\documentclass[a4,11pt]{aleph-notas}

% -- Paquetes adicionales
\usepackage{enumitem}
\usepackage{aleph-comandos} 

% -- Datos
\institucion{Facultad de Ciencias Exactas, Naturales y Ambientales}
\carrera{Catálogo STEM}
\asignatura{Álgebra lineal}
\tema{Resumen 08: Coordenadas y cambio de base}
\autor[A. Merino]{Andrés Merino}
\fecha{Periodo 2026-1} 

\logouno[0.16\textwidth]{Logos/logoPUCE_04_ac}
\definecolor{colortext}{HTML}{1957d1}
\definecolor{colordef}{HTML}{1957d1}
\fuente{montserrat}

\begin{document}

\encabezado

%%%%%%%%%%%%%%%%%%%%%%%%%%%%%%%%%%%%%%
\section{Componentes y Cambio de base}
%%%%%%%%%%%%%%%%%%%%%%%%%%%%%%%%%%%%%%

En esta sección, asumiremos $n\in\N^*$.

\begin{defi}[Base ordenada]
    Sean $(E,+,\cdot,\K)$ un espacio vectorial tal que $\dim(E)=n$ y $B\subseteq E$ una base de $E$. Si fijamos un orden en los vectores de $B$, obteniendo
    \[
        B=(v_1,v_2,\ldots,v_n),
    \]
    se dice que $B$ es una base ordenada.
\end{defi}

\begin{defi}[Vector de componentes]
    Sean $(E,+,\cdot,\K)$ un espacio vectorial y 
    \[
        B=(v_1,v_2,\ldots,v_n)
    \]
    una base ordenada. Para $u\in E$, se define el vector de componentes de $u$ por
    \[
        [u]_B = \begin{pmatrix}
            \alpha_1\\\alpha_2\\\vdots\\\alpha_n
        \end{pmatrix}
    \]
    donde $\alpha_1,\alpha_2,\ldots,\alpha_n\in\K$ son tales que
    \[
        u = \alpha_1 v_1+\alpha_2 v_2+\cdots+\alpha_n v_n.
    \]
\end{defi}

\begin{advertencia}
    En la literatura, también se puede encontrar la definición anterior bajo el nombre de vector de coordenadas.
\end{advertencia}

\begin{defi}[Matriz de cambio de base]
    Sean $(E,+,\cdot,\K)$ un espacio vectorial tal que $\dim(E)=n$ y $S,T$ bases ordenadas de $E$. Se llama matriz de cambio de base de $T$ a $S$ a la única matriz de $\Mat[\K]{n}{n}$, denotada por $P_{S\gets T}$, que cumple que
    \[
        [v]_S = P_{S\gets T} [v]_T
    \]
    para todo $v\in E$.
\end{defi}

\begin{advertencia}
    En la literatura, también se puede encontrar la definición anterior bajo el nombre de matriz de transición.
\end{advertencia}


\begin{teo}
    Sean $(E,+,\cdot,\K)$ un espacio vectorial tal que $\dim(E)=n$, $S,T$ bases ordenadas de $E$ y $P_{S\gets T}$ la matriz de cambio de base de $T$ a $S$. Se tiene que $P_{S\gets T}$ es no singular y
    \[
        (P_{S\gets T})^{-1} = P_{T\gets S}.
    \]
\end{teo}

\end{document}
